\chapter{AI Component}
\label{chap:ai-component}

\section{Business Context \& AI Integration}
\label{section:business-context-and-AI-integration}

\textbf{Objective:}

\begin{figure}[H]
    \centering
    \includegraphics[width=0.8\textwidth]{chapter_5/Domain_Diagram.png}
    \caption{System's Domain Diagram}
\end{figure}

The domain diagram above illustrates the overall system structure of Vid-Scratch. The \textit{Poisoner} component, highlighted in the diagram, represents the part of the system where AI is integrated to generate adversarial perturbations and apply them to the frames of the input video.

AI is suitable for this task because generating effective adversarial perturbations requires understanding how a model processes and classifies video input, particularly the feature representations and decision boundaries on which the model relies. Rather than applying random noise, the adversarial generation process uses gradient-based optimization to produce perturbations that are specifically designed to affect the model's predictions while remaining imperceptible to human viewers.

By applying perturbations that exploit the model's internal feature representations, the system aims to cause I3D to incorrectly classify the action shown in the video at inference time, or to learn distorted patterns when trained on poisoned data. At the same time, the visual appearance of the output video remains close to the original, allowing the protected content to remain usable for human viewers.

\section{Goal Hierarchy}
\label{section:goal-hierarchy}

\subsubsection{Organization Goals}
\begin{enumerate}
    \item 
    \textbf{Goal:}
    Develop an accessible application for generating adversarial video inputs targeting I3D within the defined project scope. \\
    \textbf{Measured by:} 
    Successful delivery of a working application that supports the generation of adversarial video outputs for both inference-time and training-time evaluation scenarios. \\
    \textbf{Data collection:}
    System functionality testing and academic evaluation by project supervisors. \\
    \textbf{Operationalization:} 
    Verify that the application can generate adversarial video outputs and that the produced results are consistent with the expected behavior observed in both evaluation scenarios.
\end{enumerate}

\subsubsection{System Goals}
\begin{enumerate}
    \item 
    \textbf{Goal:}
    Generate adversarial video outputs that influence I3D model behavior while preserving visual similarity to the original video. \\
    \textbf{Measured by:} 
    Attack success rate during inference-time evaluation, performance degradation during training-time evaluation, and perceptual similarity measurements such as SSIM. \\
    \textbf{Data collection:}
    Model prediction logs, classification accuracy records, and perceptual similarity measurements on processed video samples. \\
    \textbf{Operationalization:} 
    Evaluate whether adversarially perturbed videos cause incorrect action classification by I3D, and whether poisoned training data reduces model accuracy, while confirming that SSIM remains above an acceptable threshold.
\end{enumerate}

\subsubsection{User Goals}
\begin{enumerate}
    \item 
    \textbf{Goal:}
    Allow users to generate adversarial video variants through a simple and accessible workflow without requiring deep technical knowledge. \\
    \textbf{Measured by:} 
    User satisfaction ratings and feedback on the usability and clarity of the application. \\
    \textbf{Data collection:}
    Post-use feedback forms and usability observations during system testing. \\
    \textbf{Operationalization:} 
    Collect and analyze user feedback to assess whether the workflow is clear, easy to use, and produces outputs that meet expectations in terms of visual similarity.
\end{enumerate}

\subsubsection{AI Model Goals}
\begin{enumerate}
    \item 
    \textbf{Goal:}
    Generate adversarial perturbations that are effective against I3D by exploiting both the spatial and temporal structure of the input video. \\
    \textbf{Measured by:} 
    Attack success rate during inference-time evaluation and accuracy reduction during training-time evaluation, compared with a baseline of unperturbed inputs. \\
    \textbf{Data collection:}
    I3D prediction results on adversarial and clean video inputs, together with training accuracy curves on poisoned and clean datasets. \\
    \textbf{Operationalization:} 
    Compare model predictions and training outcomes between clean and adversarial conditions to quantify the impact of the generated perturbations on I3D performance.
\end{enumerate}

\section{Task Requirements Analysis Using AI Canvas}
\label{section:task-requirements-analysis-using-AI-canvas}

\subsection{AI Task Requirements}

\subsubsection*{Requirements (REQ)}
\begin{itemize}
    \item REQ-1: The AI must generate adversarial perturbations 
    that cause a pre-trained I3D model to incorrectly classify 
    the action shown in the input video at inference time.
    \item REQ-2: The AI must generate perturbed video samples 
    that reduce the classification performance of an I3D model 
    trained on the perturbed dataset.
    \item REQ-3: The generated perturbations must preserve 
    visual similarity to the original video as measured by 
    perceptual similarity metrics such as SSIM.
    \item REQ-4: The AI must identify important frames for 
    perturbation in order to improve attack effectiveness while 
    keeping the number of modified frames minimal.
\end{itemize}

\subsubsection*{Specifications (SPEC)}
\begin{itemize}
    \item SPEC-1: The system implements the DeepSAVA pipeline, 
    combining Bayesian Optimization for frame selection, spatial 
    transformation via a learned flow field, and gradient-based 
    additive noise optimization.
    \item SPEC-2: The perturbation strength $\epsilon$ and SSIM 
    budget are configurable by the user through the application 
    interface.
    \item SPEC-3: Perturbations are applied to a 40-frame 
    representation of the input video sampled uniformly from 
    the full clip.
    \item SPEC-4: For the training-time attack scenario, 
    Error-Minimizing Noise is applied to non-attacked frames.
\end{itemize}

\subsubsection*{Environment (ENV)}
\begin{itemize}
    \item ENV-1: The system runs in a Python environment 
    with PyTorch as the deep learning backend.
    \item ENV-2: Input data consists of MP4 video files 
    of up to one minute in length.
    \item ENV-3: GPU acceleration is recommended for efficient 
    adversarial generation, particularly for the optimization 
    steps in the DeepSAVA pipeline.
    \item ENV-4: The backend is served via a FastAPI server 
    and accessed through an Electron-based frontend application.
\end{itemize}

\subsection{AI Canvas Development}

\begin{table}[H]
\centering
\begin{tabular}{|p{4cm}|p{10cm}|}
\hline
\textbf{Canvas Element} & \textbf{Description} \\
\hline
\textbf{Input Data} & MP4 video files of up to one minute 
in length, uploaded by the user through the application 
interface. \\
\hline
\textbf{Expected Output} & Adversarially perturbed video 
files that cause I3D to incorrectly classify the action 
at inference time, or perturbed video samples that reduce 
model performance when used as training data, while 
remaining visually similar to the original. \\
\hline
\textbf{Tools/Frameworks} & Python, PyTorch, FastAPI, 
OpenCV, and Electron. The adversarial generation pipeline 
is based on DeepSAVA, incorporating Bayesian Optimization, 
spatial transformation, and gradient-based noise 
optimization. \\
\hline
\textbf{Success Criteria} & Attack success against I3D is 
observed during inference-time evaluation; training on 
perturbed data results in measurable performance reduction 
compared with clean data; and SSIM between the original 
and output video remains above an acceptable threshold. \\
\hline
\end{tabular}
\end{table}

\section{User Experience Design with AI}

\label{section:user-experience-design-with-AI}

\subsection{Interaction Style}
\textbf{Interaction Style: Process and Display}

The user interacts with the AI component indirectly through 
the application interface. The user selects an input video, 
configures generation parameters, and starts the processing 
pipeline. The AI then generates adversarial perturbations in 
the background and returns the output video for the user to 
preview and export. The system displays both the original 
and output videos side by side so that the user can verify 
that visual similarity has been preserved.

\subsection{User Feedback Mechanism}
Users are encouraged to provide feedback to help improve 
the system. The application offers a dedicated Feedback 
page with the following options:
\begin{itemize}
    \item \textbf{Comments and Suggestions} — Share ideas 
    for improving the adversarial generation pipeline or 
    the application interface.
    \item \textbf{Bug Reports} — Report technical issues or 
    unexpected behavior during video processing.
    \item \textbf{Email Us} — Contact the development team 
    directly for additional support.
\end{itemize}

\subsection{AI Contribution to System Intelligence}

Vid-Scratch integrates AI to perform gradient-based 
adversarial video generation targeting I3D. This approach 
provides capabilities that would be difficult to achieve 
through non-AI methods alone.

\subsubsection{Without AI Integration}
\begin{itemize}
    \item Manual frame editing or fixed noise filters that 
    do not account for how the model processes video input.
    \item Static perturbations that may be less effective 
    against temporal aggregation in video models such as I3D.
\end{itemize}

\subsubsection{With AI Integration}
\begin{itemize}
    \item Gradient-based optimization that targets the 
    feature representations used by I3D.
    \item Bayesian Optimization for identifying the most 
    influential frames to perturb, improving perturbation 
    effectiveness while keeping modifications minimal.
    \item Spatial transformation combined with additive 
    noise to exploit both the spatial and temporal structure 
    of the video.
    \item Preservation of visual similarity between the 
    original and output video through SSIM-constrained 
    optimization.
\end{itemize}
